\chapter{NGUỒN DỮ LIỆU VÀ PHƯƠNG THỨC THU THẬP}\label{chap:dulieu}

Dữ liệu là bài báo khoa học trên arXiv, thu thập qua ba kênh: arXiv API (theo keyword, chuyên mục, khoảng ngày
hoặc danh sách mã bài), bộ dữ liệu arXiv trên Kaggle, và PDF do người dùng tải lên. Tải file PDF là một tùy chọn
của kênh arXiv API, không phải kênh riêng. Bộ dữ liệu thử nghiệm gồm 492 bài, trong đó 12 bài có toàn văn PDF;
ngoài ra có 4 PDF được giữ riêng để trình diễn thao tác thêm bài.

\section{Các kênh thu thập}

\begin{table}[H]
\centering\small\setstretch{1.15}
\begin{tabularx}{\linewidth}{@{}P{0.2\linewidth}LP{0.17\linewidth}P{0.24\linewidth}@{}}
\toprule
\textbf{Kênh} & \textbf{Nội dung lấy được} & \textbf{Quy mô phù hợp} & \textbf{Lệnh}\\
\midrule
arXiv API~\cite{arxivapi} (Atom feed) & Mã bài, tiêu đề, tóm tắt, tác giả, chuyên mục, ngày, DOI, link PDF; tùy chọn tải PDF (toàn văn, mục, tài liệu tham khảo) & Vài nghìn bài mỗi truy vấn & \code{crawler keyword}, \code{crawler category}, \code{crawler ids}, cờ \code{--pdf}\\
Kaggle dataset Cornell-University/arxiv~\cite{arxivdataset} & Metadata của khoảng 2,7 triệu bài (JSON lines) & Hàng trăm nghìn bài & \code{crawler snapshot}\\
Người dùng tải lên & PDF bất kỳ, có hoặc không có mã arXiv; hệ thống tự tra metadata chuẩn trên arXiv API & Từng file hoặc theo lô & Giao diện web, \code{ingest pdf}, \code{ingest dir}\\
\bottomrule
\end{tabularx}
\caption{Các kênh thu thập dữ liệu}\label{tab:kenh}
\end{table}

\section{Phương thức thu thập}

\begin{enumerate}[label=(\arabic*)]
\item \textbf{Dựng truy vấn:} keyword được chuyển sang cú pháp arXiv. Cụm trong ngoặc kép thành
  \code{all:"..."}, từ rời được nối bằng AND; có thể ghép thêm chuyên mục \code{cat:cs.IR} và khoảng ngày
  \code{submittedDate:[từ TO đến]}.
\item \textbf{Thu thập không theo keyword:} lấy theo chuyên mục và khoảng ngày, ví dụ
  \code{crawler category cs.IR --from 2025-01-01}.
\item \textbf{Phân trang và tuân thủ quy định:} mỗi trang 100 kết quả, sắp xếp mới nhất trước; nghỉ tối thiểu
  3 giây giữa các lần gọi theo quy định của arXiv; thử lại khi lỗi mạng hoặc HTTP khác 200.
\item \textbf{Chống trùng khi thu thập:} kết quả được ghi nối vào \code{data/raw/metadata.jsonl}; mã bài đã có bị
  bỏ qua, nên thu thập lại nhiều lần không sinh bản sao.
\item \textbf{Tải PDF:} kiểm tra chữ ký \code{\%PDF} và kích thước file trước khi lưu vào \code{data/pdfs/}.
\end{enumerate}

Sau bước thu thập, bài báo chưa xuất hiện trong ứng dụng; cần chạy tiếp các bước nạp (\code{ingest}), embedding và
dựng đồ thị (Chương~\ref{chap:tochuc}).

\section{Bộ dữ liệu thử nghiệm}

\begin{table}[H]
\centering\small\setstretch{1.15}
\begin{tabularx}{\linewidth}{@{}Lr@{}}
\toprule
\textbf{Truy vấn thu thập} & \textbf{Số bài mới}\\
\midrule
12 bài kinh điển theo mã bài, kèm PDF (Transformer, BERT, RAG, DPR, REALM, FiD, Contriever, E5, BGE, Sentence-BERT, Self-RAG, RAPTOR) & 12\\
Keyword ``retrieval augmented generation'' & 150\\
Keyword ``dense retrieval'' & 97\\
Keyword ``graph neural network'', chuyên mục cs.LG & 80\\
Keyword ``knowledge graph'' question answering & 76\\
Chuyên mục cs.IR từ 01/2025 & 77\\
\midrule
\textbf{Tổng} & \textbf{492}\\
\bottomrule
\end{tabularx}
\caption{Thành phần bộ dữ liệu thử nghiệm}\label{tab:dataset}
\end{table}

Bốn PDF dành cho demo (không nạp sẵn): GraphRAG~\cite{graphrag} (2404.16130), RAG Survey (2312.10997),
LightRAG~\cite{lightrag} (2410.05779) và Agentic RAG Survey (2501.09136). Các file này được tải lại bằng
\code{scripts/download\_demo\_pdfs.py}.
